\documentclass[10pt,a4paper]{amsart}
\usepackage[margin=19mm]{geometry}
\usepackage{amsmath,amssymb,mathtools}
\usepackage[scr=rsfs,cal=euler]{mathalfa}
\usepackage{xfrac}
\usepackage[unicode,hidelinks,bookmarks=false]{hyperref}
\newcommand{\ie}{i.e.\ }
\newcommand{\cf}{cf.\ }
\newcommand{\resp}{respectively\ }
\newcommand{\Subsection}{\subsection}
\providecommand{\nobreakdash}{\nobreak\hbox{-}}
\setlength{\emergencystretch}{2em}
\multlinegap0pt
\usepackage{bm}
\usepackage{bbm}
\usepackage{dsfont}
\usepackage{xspace}
\usepackage{cancel}
\usepackage{mathtools}
\usepackage{amsthm}
\makeatletter
\@ifundefined{theorem}{\newtheorem{theorem}{Theorem}}{}
\@ifundefined{lemma}{\newtheorem{lemma}[theorem]{Lemma}}{}
\@ifundefined{proposition}{\newtheorem{proposition}[theorem]{Proposition}}{}
\makeatother
\newcommand{\ind}{\mathbf 1}
\title[Residues of a tropical zeta function for convex domains]{Residues of a tropical zeta function for convex domains}
\author{Nikita Kalinin, Ernesto Lupercio,, Mikhail Shkolnikov}
\date{Source: arXiv:2604.21709v2 (CC BY 4.0)}
\begin{document}
\maketitle

\noindent\emph{Source and licence.} Residues of a tropical zeta function for convex domains. Version arXiv:2604.21709v2 (CC BY 4.0), distributed under the Creative Commons Attribution 4.0 International licence, as recorded in the arXiv licence metadata for this exact version and shown on the abstract page https://arxiv.org/abs/2604.21709v2. Full rights notice: licenses/arxiv-2604.21709-CC-BY-4.0.txt.\par\smallskip

\noindent\textit{Editorial scope.} Counted unit: the Proposition whose source label is \texttt{prop:spec-strip} in the article (source0.tex lines 7436--7465, source label \texttt{prop:spec-strip}), delivered together with the article's own complete original proof of it (source0.tex lines 7467--7746, one \texttt{proof} environment of 280 lines). No mathematics is written, repaired or re-derived: statement and proof are the article's own text line for line. Required context retained uncounted: eq:ass (lines 1722--1728); prop:spec (lines 2002--2021); lem:unitavg (lines 6505--6515); lem:inckloo (lines 6630--6644); lem:fejer\_moment (lines 6793--6803); lem:Hs-strip (lines 7250--7268); lem:Bs-strip (lines 7349--7372). Every definition, display and label of those lines is retained unchanged. Omitted from this package and not redistributed: 1--1721, 1729--2001, 2022--6504, 6516--6629, 6645--6792, 6804--7249, 7269--7348, 7373--7435, 7466--7466, 7747--8987. Nothing inside a retained excerpt is omitted or altered: every retained body is delivered exactly as the article's lines are written, so no omission receipt of any kind accompanies this package. Citations: the retained excerpts carry no citation command at all, so no bibliography entry of the article is transcribed and no reference is redistributed. Reference adaptation: none was needed and none was made; every reference inside the delivered unit resolves inside it, so no unresolved reference or citation is produced. Printed slips kept literal: no printed word, symbol, hypothesis or number of the counted statement or of its proof was altered. Layout and class: the article's own class is \texttt{amsart} with packages including babel, csquotes, amsmath, amssymb, amsfonts, amscd, graphicx, caption, subcaption, wrapfig, pdflscape, float, listings, xcolor; the wrapper is the lane's pinned \texttt{amsart} editorial wrapper at 10pt on A4, which restates the theorem-like environments the retained text needs and adds the article's own macro definitions that the retained text needs (\texttt{\textbackslash ind}), with extra wrapper packages (bm, bbm, dsfont, xspace, cancel, mathtools). The delivered numbering is the unit's own. Assets: no retained span contains a figure environment, a graphics-inclusion command, a PDF-page inclusion, a table or a TikZ picture; no image file is copied into this package, the package declares no asset dependency and no image file is redistributed with it. Licence: version arXiv:2604.21709v2 of this article is distributed under the Creative Commons Attribution 4.0 International licence, as recorded in the arXiv licence metadata for this exact version and shown on the abstract page https://arxiv.org/abs/2604.21709v2. A full rights notice is delivered at licenses/arxiv-2604.21709-CC-BY-4.0.txt. Characters: every retained excerpt is pure ASCII, so the delivered engine, XeLaTeX with its default Latin Modern text font, reports no missing character.
\begin{equation}
\label{eq:ass}
f\in C^3([0,1])
\qquad\text{and}\qquad
0<m\le |f''(x)|
\quad\text{for all }x\in[0,1].
\end{equation}

\begin{proposition}
\label{prop:spec}
Fix a compact set
\[
K\subset\left\{\frac12<\Re(s)<1\right\}
\]
and assume \eqref{eq:ass}. Then for every \(\varepsilon>0\) and \(s\in K\)
we have
\[
\Sigma_b(s)
=
\varphi(b)
\left(\int_0^1 H_s(u)\,du\right)
\left(\int_0^1 |f''(v)|^s\,dv\right)
+
O_{f,K,\varepsilon}
\bigl(b^{\frac{1+\Re(s)}{2}+\varepsilon}\bigr),
\]
with implied constant uniform for \(s\in K\).
\end{proposition}

\begin{lemma}[Reduced residues: $H_s$ averages to its integral]
\label{lem:unitavg}
Fix a compact set $K\Subset\{\tfrac12<\Re(s)<1\}$. Then for $s\in K$ and $b\ge1$,
\[
\sum_{\substack{1\le r\le b\\(r,b)=1}} H_s(r/b)
=
\varphi(b)\int_0^1 H_s(u)\,du
+O_K\bigl(b^{\Re(s)}\tau(b)\bigr),
\]
where $\tau(b)$ is the divisor function and the implied constant is uniform for $s\in K$.
\end{lemma}

\begin{lemma}[Completion bound for \(K_b(n;R)\)]
\label{lem:inckloo}
For all integers \(n\), all \(b\ge 1\), and all \(1\le R\le b\),
\[
K_b(n;R)
\ll
(n,b)^{1/2}\,b^{1/2}\,\tau(b)^2\log(2b),
\]
with an absolute implied constant.

In particular, if \((n,b)=1\), then for every \(\varepsilon>0\),
\[
\max_{1\le R\le b}|K_b(n;R)|\ll_{\varepsilon} b^{1/2+\varepsilon}.
\]
\end{lemma}

\begin{lemma}[First moment of the Fej\'er kernel]
\label{lem:fejer_moment}
There exists an absolute constant \(C>0\) such that for all \(N\ge2\),
\[
\int_{-1/2}^{1/2} |t|\,F_N(t)\,dt\le C\,\frac{\log N}{N}.
\]
Consequently, if \(G\) is \(1\)-periodic and Lipschitz, then its Fej\'er mean \(G*F_N\) satisfies
\[
\|G-G*F_N\|_{\infty}\le C\,\frac{\log N}{N}\,\operatorname{Lip}(G).
\]
\end{lemma}

\begin{lemma}[Strip-uniform bounds for \(H_s\) and \(\partial_u H_s\)]
\label{lem:Hs-strip}
Fix numbers
\[
\frac12<\alpha<\beta<1.
\]
Then there exist constants \(C_1=C_1(\alpha)\) and \(C_2=C_2(\alpha,\beta)\) such that for every
\[
s=\sigma+i\tau,\qquad \alpha\le \sigma\le \beta,
\]
and every \(u\in(0,1]\), one has
\[
|H_s(u)|\le C_1\,u^{-\sigma},
\]
and
\[
|\partial_u H_s(u)|\le C_2\,(1+|\tau|)\,u^{-\sigma-1}.
\]
\end{lemma}

\begin{lemma}[Strip-uniform bounds for \(B_s(x)=|f''(x)|^s\)]
\label{lem:Bs-strip}
Fix numbers
\[
\frac12<\alpha<\beta<1.
\]
Assume
\[
f\in C^3([0,1]),
\qquad
0<m\le |f''(x)|\le M
\quad (x\in[0,1]).
\]
Then there exist constants \(C_0,C_1,C_2>0\), depending only on \(\alpha,\beta,f\), such that for all
\[
s=\sigma+i\tau,\qquad \alpha\le \sigma\le \beta,
\]
and all \(x\in[0,1]\), one has
\[
|B_s(x)|\le C_0,\qquad \|B_s\|_\infty\le C_1,\qquad \|B_s'\|_\infty\le C_2(1+|\tau|).
\]
In particular, the affine interpolant/remainder decomposition used in the proof of
Proposition~\ref{prop:spec} remains valid on strips, with polynomial dependence on \(|\tau|\).
\end{lemma}

\begin{proposition}[Strip-uniform equidistribution estimate]
\label{prop:spec-strip}
Fix numbers
\[
\frac12<\alpha<\beta<1.
\]
Assume
\[
f\in C^3([0,1]),
\qquad
0<m\le |f''(x)|\le M
\quad (x\in[0,1]).
\]
Then for every \(\varepsilon>0\) there exists \(C=C(\alpha,\beta,f,\varepsilon)>0\) such that for every
\[
s=\sigma+i\tau,\qquad \alpha\le \sigma\le \beta,
\]
and every \(b\ge1\), one has
\[
\Sigma_b(s)
=
\varphi(b)\left(\int_0^1 H_s(u)\,du\right)\left(\int_0^1 |f''(v)|^s\,dv\right)
+
O\!\left((1+|\tau|)^2\,b^{\frac{1+\sigma}{2}+\varepsilon}\right),
\]
where
\[
\Sigma_b(s):=\sum_{\substack{1\le r\le b\\(r,b)=1}} H_s(r/b)\,|f''(\overline r/b)|^s.
\]
\end{proposition}

\begin{proof}
Write
\[
B_s(x):=|f''(x)|^s,
\qquad
A_r:=H_s(r/b).
\]

We first record the strip-uniform analogue of Lemma~\ref{lem:unitavg}:
\begin{equation}
\label{eq:unitavg-strip}
\sum_{\substack{1\le r\le b\\(r,b)=1}} H_s(r/b)
=
\varphi(b)\int_0^1 H_s(u)\,du
+
O_{\alpha,\beta}\!\left((1+|\tau|)\,b^\sigma \tau(b)\right).
\end{equation}
Indeed, the proof of Lemma~\ref{lem:unitavg} carries over verbatim once one replaces the
compact-uniform bounds used there by the strip-uniform bounds of Lemma~\ref{lem:Hs-strip}.
The only place where \(|\tau|\) appears is through the derivative bound for \(H_s\), and this yields
the factor \((1+|\tau|)\).

Now follow the proof of Proposition~\ref{prop:spec}. Let
\[
\ell_s(x):=(1-x)B_s(0)+xB_s(1),
\qquad
G_s(x):=B_s(x)-\ell_s(x).
\]
By Lemma~\ref{lem:Bs-strip},
\[
\|B_s\|_\infty\ll_{\alpha,\beta,f}1,
\qquad
\operatorname{Lip}(G_s)\ll_{\alpha,\beta,f}(1+|\tau|).
\]

Set
\[
N:=R:=\lfloor b^{1/2}\rfloor.
\]
By the Fej\'er-kernel estimate (Lemma~\ref{lem:fejer_moment}),
\[
\|G_s-(G_s)_N\|_\infty
\ll
\frac{\log(2N)}{N}\operatorname{Lip}(G_s)
\ll_{\alpha,\beta,f}
(1+|\tau|)\frac{\log(2N)}{N}.
\]
Therefore
\[\mkern-25mu
\Sigma_b(s)
=
\sum_{(r,b)=1}A_r\,(G_s)_N(\overline r/b)
+
\sum_{(r,b)=1}A_r\,\ell_s(\overline r/b)
+
O\!\left((1+|\tau|)\frac{\log(2N)}{N}\sum_{(r,b)=1}|A_r|\right).
\]
By Lemma~\ref{lem:Hs-strip},
\[
|A_r|=|H_s(r/b)|\ll_\alpha (r/b)^{-\sigma},
\]
so
\[
\sum_{(r,b)=1}|A_r|
\le
\sum_{r=1}^b |H_s(r/b)|
\ll_\alpha
b^\sigma\sum_{r=1}^b r^{-\sigma}
\ll_\alpha b.
\]
Hence the truncation error is
\begin{equation}
\label{eq:trunc-strip}
O\!\left((1+|\tau|)\,b^{1/2}\log(2b)\right).
\end{equation}

Next expand the periodic part:
\[
(G_s)_N(x)=\widehat G_s(0)+\sum_{0<|n|<N}\widehat{(G_s)_N}(n)e(nx).
\]
Since \(G_s\) is Lipschitz with
\[
\operatorname{Lip}(G_s)\ll (1+|\tau|),
\]
its Fourier coefficients satisfy
\[
|\widehat G_s(n)|\ll_{\alpha,\beta,f}\frac{1+|\tau|}{|n|}
\qquad (n\neq0),
\]
hence
\[
|\widehat{(G_s)_N}(n)|\le |\widehat G_s(n)|
\ll_{\alpha,\beta,f}\frac{1+|\tau|}{|n|}.
\]
So
\[
\sum_{(r,b)=1}A_r\,(G_s)_N(\overline r/b)
=
\widehat G_s(0)\sum_{(r,b)=1}A_r
+
\sum_{0<|n|<N}\widehat{(G_s)_N}(n)\,S_b(n),
\]
where
\[
S_b(n):=\sum_{(r,b)=1}A_r\,e\!\left(\frac{n\overline r}{b}\right).
\]

For the affine part, write
\[
\ell_s(x)=\alpha_s+\beta_s x,
\qquad \beta_s=B_s(1)-B_s(0).
\]
Since \(\|B_s\|_\infty\ll1\), we have \(|\beta_s|\ll1\). The discrete Fourier expansion on the \(b\)-grid gives
\[
\ell_s(m/b)=c_{b,s}(0)+\sum_{0<|n|\le b/2}c_{b,s}(n)e(nm/b),
\]
with
\[
c_{b,s}(0)=\int_0^1 \ell_s(x)\,dx+O(b^{-1}),
\qquad
|c_{b,s}(n)|\ll_{\alpha,\beta,f}\frac{1}{|n|}
\quad (0<|n|\le b/2).
\]
Therefore
\[
\sum_{(r,b)=1}A_r\,\ell_s(\overline r/b)
=
c_{b,s}(0)\sum_{(r,b)=1}A_r
+
\sum_{0<|n|\le b/2}c_{b,s}(n)S_b(n).
\]

We now bound \(S_b(n)\). Split
\[
S_b(n)=S_b^{\le R}(n)+S_b^{>R}(n),
\]
where
\[
S_b^{\le R}(n):=\sum_{\substack{1\le r\le R\\(r,b)=1}}A_r e\!\left(\frac{n\overline r}{b}\right),
\qquad
S_b^{>R}(n):=\sum_{\substack{R<r\le b\\(r,b)=1}}A_r e\!\left(\frac{n\overline r}{b}\right).
\]

For the initial segment,
\[
|S_b^{\le R}(n)|
\le
\sum_{r\le R}|A_r|
\ll_\alpha
b^\sigma\sum_{r\le R}r^{-\sigma}
\ll_\alpha
b^\sigma R^{1-\sigma}
\ll_\alpha
b^{\frac{1+\sigma}{2}}.
\]

For the tail, define
\[
a_r:=\ind_{(r,b)=1}e\!\left(\frac{n\overline r}{b}\right),
\qquad
T_n(T):=\sum_{1\le r\le T}a_r.
\]
By Lemma~\ref{lem:inckloo},
\[
\max_{1\le T\le b}|T_n(T)|
\ll_\eta
(n,b)^{1/2}b^{1/2+\eta}
\]
for any \(\eta>0\). Abel summation gives
\[
S_b^{>R}(n)
=
H_s(1)\,T_n(b)-H_s(R/b)\,T_n(R)
-\sum_{T=R}^{b-1}T_n(T)\Bigl(H_s((T+1)/b)-H_s(T/b)\Bigr).
\]
By Lemma~\ref{lem:Hs-strip},
\[
|H_s(1)|\ll_\alpha 1,
\qquad
|H_s(R/b)|\ll_\alpha (R/b)^{-\sigma}.
\]
Also, by the mean value theorem and Lemma~\ref{lem:Hs-strip},
\[
|H_s((T+1)/b)-H_s(T/b)|
\le
\frac1b\sup_{u\in[T/b,(T+1)/b]}|\partial_u H_s(u)|
\ll_{\alpha,\beta}
(1+|\tau|)\,b^\sigma T^{-\sigma-1}.
\]
Hence
\[
\sum_{T=R}^{b-1}|H_s((T+1)/b)-H_s(T/b)|
\ll_{\alpha,\beta}
(1+|\tau|)\,b^\sigma\sum_{T\ge R}T^{-\sigma-1}
\ll_{\alpha,\beta}
(1+|\tau|)\,b^\sigma R^{-\sigma}.
\]
Therefore
\[
|S_b^{>R}(n)|
\ll_{\alpha,\beta,\eta}
(n,b)^{1/2}b^{1/2+\eta}
\left(
1+(R/b)^{-\sigma}+(1+|\tau|)b^\sigma R^{-\sigma}
\right).
\]
Since \(R=\lfloor b^{1/2}\rfloor\), both \((R/b)^{-\sigma}\) and \(b^\sigma R^{-\sigma}\) are
\(\ll b^{\sigma/2}\). Thus
\[
S_b(n)\ll_{\alpha,\beta,\eta}
(1+|\tau|)\,(n,b)^{1/2}\,b^{\frac{1+\sigma}{2}+\eta}.
\]

Now sum the nonzero Fourier modes. For the periodic part,
\[
\sum_{0<|n|<N}|\widehat{(G_s)_N}(n)S_b(n)|
\ll_{\alpha,\beta,f,\eta}
(1+|\tau|)^2 b^{\frac{1+\sigma}{2}+\eta}
\sum_{1\le n<N}\frac{(n,b)^{1/2}}{n}.
\]
Grouping by \(d=(n,b)\) exactly as in the proof of Proposition~\ref{prop:spec},
\[
\sum_{1\le n<N}\frac{(n,b)^{1/2}}{n}
\ll \tau(b)\log(2N).
\]
Therefore
\[
\sum_{0<|n|<N}\widehat{(G_s)_N}(n)S_b(n)
=
O\!\left((1+|\tau|)^2 b^{\frac{1+\sigma}{2}+\eta}\tau(b)\log(2N)\right).
\]

Similarly, for the affine part,
\[
\sum_{0<|n|\le b/2}|c_{b,s}(n)S_b(n)|
\ll_{\alpha,\beta,f,\eta}
(1+|\tau|)\,b^{\frac{1+\sigma}{2}+\eta}
\sum_{1\le n\le b/2}\frac{(n,b)^{1/2}}{n},
\]
hence
\[
\sum_{0<|n|\le b/2}c_{b,s}(n)S_b(n)
=
O\!\left((1+|\tau|)\,b^{\frac{1+\sigma}{2}+\eta}\tau(b)\log(2b)\right).
\]

For the zero mode,
\[
\widehat G_s(0)+c_{b,s}(0)
=
\int_0^1 B_s(x)\,dx+O(b^{-1})
=
\int_0^1 |f''(x)|^s\,dx+O(b^{-1}),
\]
so by \eqref{eq:unitavg-strip},
\[
\bigl(\widehat G_s(0)+c_{b,s}(0)\bigr)\sum_{(r,b)=1}A_r
=
\varphi(b)\left(\int_0^1 H_s(u)\,du\right)\left(\int_0^1 |f''(x)|^s\,dx\right)
+
O\!\left((1+|\tau|)\,b^\sigma\tau(b)\right)
+
O(1).
\]

Since \(\sigma<1\), we have
\[
b^\sigma\tau(b)\ll_\varepsilon b^{\frac{1+\sigma}{2}+\varepsilon},
\]
and because \(\sigma>\tfrac12\),
\[
b^{1/2}\log(2b)\ll_\varepsilon b^{\frac{1+\sigma}{2}+\varepsilon}.
\]
Absorbing the factors \(\tau(b)\log(2b)\) into \(b^\varepsilon\), and then choosing \(\eta>0\)
sufficiently small relative to the final \(\varepsilon\), all the error terms combine into
\[
O\!\left((1+|\tau|)^2\,b^{\frac{1+\sigma}{2}+\varepsilon}\right).
\]
This proves the proposition.
\end{proof}

\end{document}
